% sdapsclassic: SDAPS document class for questionnaires.
% Options:
%   14pt         — base font size; 14pt is recommended for older adults
%   english      — language used for SDAPS-generated text (e.g. instructions)
%   a4paper      — page size
%   sdaps_style=code128 — barcode style printed at the bottom of each page;
%                  code128 is a 1-D barcode that encodes page/questionnaire ID
%   checkmode=check — how marks are detected: "check" accepts any mark in the box
\documentclass[
  14pt,
  english,
  a4paper,
  sdaps_style=code128,
  checkmode=check,
  print_questionnaire_id
]{sdapsclassic}

% Charter: a serif font with a large x-height and open letterforms,
% recommended for readability in print questionnaires for older adults.
\usepackage{charter}

% Override the normal font size to 14pt with 18pt leading (line spacing).
% The two numbers are {font size}{baseline skip (line spacing)}.
% Increasing baseline skip beyond 1.2× font size improves readability.
\renewcommand{\normalsize}{\fontsize{14}{18}\selectfont}

\ExplSyntaxOn

% Suppress the "DRAFT" watermark that SDAPS prints before the questionnaire
% has been registered with the sdaps command-line tool.
\bool_gset_false:N \g_sdaps_draft_bool

% Set the size and border thickness of the response circles.
%   width/height — diameter of the circle in mm (default: 3.5mm)
%   linewidth    — thickness of the circle border (default: 1bp)
% singlechoice = radio buttons (select one); multichoice = tick boxes (select many)
\sdaps_context_append:nn { singlechoice } { width=6mm, height=6mm, linewidth=1.5bp }
\sdaps_context_append:nn { multichoice }  { width=6mm, height=6mm, linewidth=1.5bp }

% Set the size of the barcode printed at the bottom of each page.
%   c_sdaps_barcode_height_dim  — height of the bars (default: 6.5mm)
%   c_sdaps_barcode_bar_width_dim — width of each individual bar (default: 0.33mm);
%                                   increasing this makes the whole barcode wider.
%                                   Do not exceed ~0.66mm or it may not fit the margin.
\dim_gset:Nn \c_sdaps_barcode_height_dim  { 10mm }
\dim_gset:Nn \c_sdaps_barcode_bar_width_dim { 0.4mm }

\ExplSyntaxOff

\title{GDS}
\author{}
\date{}

% Suppress automatic section numbering in the output.
\setcounter{secnumdepth}{0}

\begin{document}

% questionnaire: the main SDAPS environment. Every page must be inside this.
%   noinfo — suppresses the "how to fill in this form" instructions box
%            that SDAPS normally prints at the top of the first page.
\begin{questionnaire}[noinfo]

% addinfo: stores key–value metadata in the .sdaps project file.
% scale_id is used by the downstream data pipeline to identify which
% questionnaire this is when multiple scales share a single project.
\addinfo{scale_id}{gds}

\section{Staff Use Only}

    % Side-by-side layout for the ID and Week boxes.
    % Each minipage is 48% of the line width; \hfill provides the gap between them.
    \noindent
    \begin{minipage}[t]{0.48\linewidth}
        % \textbox* creates a free-text write-in box for manual entry by staff.
        %   * (star)  — the box does NOT stretch vertically to fill extra space
        %   var=...   — variable name written to the SDAPS data file
        %   {0.8cm}   — height of the write-in box
        %   {ID:}     — label printed to the left of the box
        \textbox*[var=worker_id]{0.8cm}{ID:}
    \end{minipage}\hfill
    \begin{minipage}[t]{0.48\linewidth}
        \textbox*[var=week]{0.8cm}{Week:}
    \end{minipage}

    % Side-by-side layout for the Timepoint label and its radio buttons.
    % The label sits in a 3cm minipage; the options share the remaining width.
    % [c] vertically centers both minipages relative to each other.
    \noindent
    \begin{minipage}[c]{3cm}
        Timepoint:
    \end{minipage}%
    \begin{minipage}[c]{\dimexpr\linewidth-1cm\relax}
        % optionquestion: a single-answer (radio button) question.
        %   cols=4    — arrange the choices in 4 equal columns on one row
        %   var=...   — variable name written to the SDAPS data file
        %   {}        — empty title suppresses the subsection heading above the choices
        \begin{optionquestion}[cols=4, var=timepoint]{}
            % \choiceitem defines one radio-button option.
            % The text is the label printed next to the circle.
            \choiceitem{Pre-class}
            \choiceitem{Post-class}
            \choiceitem{Baseline}
            \choiceitem{Post}
        \end{optionquestion}
    \end{minipage}

\section{Scale}
Choose the best answer for how you felt over the \underline{past week}.

\vspace{5mm}

% optiongroup: a grid of questions sharing the same set of response options.
%   colsep=5pt  — horizontal space between columns
%   rowsep=10pt — vertical space between question rows
%   {}          — empty title (no heading above the table)
\begin{optiongroup}[colsep=5pt, rowsep=10pt]{}

    % \choice defines one column header (a response option shared by all questions).
    %   val=N — the numeric value written to the data file when this option is selected
    % \parbox[t]{1.5cm}{\centering ...} forces all column headers to the same
    %   1.5cm width so that columns are equal regardless of label length.
    %   [t] top-aligns the text; \centering centers it within the box.
    %   Increase the width if labels are cut off; decrease it to tighten the layout.
    \choice[val=1]{\parbox[t]{1.5cm}{\centering Yes}}
    \choice[val=0]{\parbox[t]{1.5cm}{\centering No}}

    % \question defines one row in the grid.
    %   var=... — variable name written to the SDAPS data file for this item
    %   {text}  — the question text displayed to the respondent
    \question[var=gds01]{Are you basically satisfied with your life?}
    \question[var=gds02]{Have you dropped many of your activities and interests?}
    \question[var=gds03]{Do you feel that your life is empty?}
    \question[var=gds04]{Do you often get bored?}
    \question[var=gds05]{Are you in good spirits most of the time?}
    \question[var=gds06]{Are you afraid that something bad is going to happen to you?}
    \question[var=gds07]{Do you feel happy most of the time?}
    \question[var=gds08]{Do you often feel helpless?}
    \question[var=gds09]{Do you prefer to stay home, rather than going out and doing things?}
    \question[var=gds10]{Do you feel that you have more problems with memory than most?}
    \question[var=gds11]{Do you think is it wonderful to be alive now?}
    \question[var=gds12]{Do you feel worthless the way you are now?}
    \question[var=gds13]{Do you feel full of energy?}
    \question[var=gds14]{Do you feel that your situation is hopeless?}
    \question[var=gds15]{Do you think that most people are better off than you are?}
\end{optiongroup}
\end{questionnaire}
\end{document}
